\subsection{Procedures as Arguments}
\label{Section 1.3.1}

Consider the following three procedures.  The first computes the sum of the
integers from \code{a} through \code{b}:

\begin{codeBlock}
>>> def sum_integers(a, b):
...     if a > b:
...         return 0
...     else:
...         return a + sum_integers(a + 1, b)
\end{codeBlock}

\noindent
The second computes the sum of the cubes of the integers in the given range:

\begin{codeBlock}
>>> def sum_cubes(a, b):
...     if a > b:
...         return 0
...     else:
...         return cube(a) + sum_cubes(a + 1, b)
\end{codeBlock}

\noindent
The third computes the sum of a sequence of terms in the series

$$ \frac{1}{1\cdot 3} +  \frac{1}{5\cdot 7} + \frac{1}{9\cdot 11} + \dots, $$

\noindent
which converges to \( \pi / 8 \) (very slowly):\footnote{This series, usually
written in the equivalent form
\( \frac{\pi}{4} = 1 - \frac{1}{3} + \frac{1}{5} - \frac{1}{7} + \dots \),
is due to Leibniz.  We'll see how to use this as the basis for some
fancy numerical tricks in \link{Section 3.5.3}.}

\begin{codeBlock}
>>> def pi_sum(a, b):
...     if a > b:
...         return 0
...     else:
...         return 1.0 / (a * (a + 2)) + pi_sum(a + 4, b)
\end{codeBlock}

\noindent
These three procedures clearly share a common underlying pattern.  They are for
the most part identical, differing only in the name of the procedure, the
function of \code{a} used to compute the term to be added, and the function
that provides the next value of \code{a}.  We could generate each of the
procedures by filling in slots in the same template:

\begin{syntaxForm}
def \meta{name}(a, b):
    if a > b:
        return 0
    else:
        return \meta{term}(a) + \meta{name}(\meta{next}(a), b)
\end{syntaxForm}

\noindent
The presence of such a common pattern is strong evidence that there is a useful
abstraction waiting to be brought to the surface.  Indeed, mathematicians long
ago identified the abstraction of \newterm{summation of a series} and invented
``sigma notation,'' for example

$$ \sum\limits_{n=a}^b f(n) = f(a) + \dots + f(b), $$

\noindent
to express this concept.  The power of sigma notation is that it allows
mathematicians to deal with the concept of summation itself rather than only
with particular sums---for example, to formulate general results about sums
that are independent of the particular series being summed.

Similarly, as program designers, we would like our language to be powerful
enough so that we can write a procedure that expresses the concept of summation
itself rather than only procedures that compute particular sums.  We can do so
readily in our procedural language by taking the common template shown above
and transforming the ``slots'' into formal parameters:

\begin{codeBlock}
>>> def sum(term, a, next, b):
...     if a > b:
...         return 0
...     else:
...         return term(a) + sum(term, next(a), next, b)
\end{codeBlock}

\noindent
Notice that \code{sum} takes as its arguments the lower and upper bounds
\code{a} and \code{b} together with the procedures \code{term} and \code{next}.
We can use \code{sum} just as we would any procedure.  For example, we can use
it (along with a procedure \code{inc} that increments its argument by 1) to
define \code{sum\_cubes}:

\begin{codeBlock}
>>> def inc(n):
...     return n + 1

>>> def sum_cubes(a, b):
...     return sum(cube, a, inc, b)
\end{codeBlock}

\noindent
Using this, we can compute the sum of the cubes of the integers from 1 to 10:

\begin{codeBlock}
>>> sum_cubes(1, 10)
~\textit{3025}~
\end{codeBlock}

\noindent
With the aid of an identity procedure to compute the term, we can define
\code{sum\_integers} in terms of \code{sum}:

\begin{codeBlock}
>>> def identity(x):
...     return x

>>> def sum_integers(a, b):
...     return sum(identity, a, inc, b)
\end{codeBlock}

\noindent
Then we can add up the integers from 1 to 10:

\begin{codeBlock}
>>> sum_integers(1, 10)
~\textit{55}~
\end{codeBlock}

\noindent
We can also define \code{pi\_sum} in the same way:\footnote{Notice that we have
used block structure (\link{Section 1.1.8}) to embed the definitions of
\code{pi\_next} and \code{pi\_term} within \code{pi\_sum}, since these
procedures are unlikely to be useful for any other purpose.  We will see how to get rid of
them altogether in \link{Section 1.3.2}.}

\begin{codeBlock}
>>> def pi_sum(a, b):
...     def pi_term(x):
...         return 1.0 / (x * (x + 2))
...     def pi_next(x):
...         return x + 4
...     return sum(pi_term, a, pi_next, b)
\end{codeBlock}

\noindent
Using these procedures, we can compute an approximation to \( \pi \):

\begin{codeBlock}
>>> 8 * pi_sum(1, 1000)
~\textit{3.139592655589783}~
\end{codeBlock}

\begin{revealTheSurface}
Python has \code{sum} already, and it is not the one we just defined---ours
shadows it for the rest of this section.  The built-in takes a sequence and
adds it up, so the two computations above are written

\begin{codeBlock}
>>> import builtins
>>> builtins.sum(cube(x) for x in range(1, 11))
~\textit{3025}~
>>> builtins.sum(range(1, 11))
~\textit{55}~
\end{codeBlock}

\noindent
The parenthesised expression is a \newterm{generator}, which produces the
terms one at a time as \code{sum} asks for them.  Between it and \code{range},
the two things our \code{sum} takes as procedures---what the term is, and what
comes next---are supplied by the sequence itself, and the procedure arguments
disappear.  For products there is \code{math.prod}, and for the general
accumulation of \link{Exercise 1.32} there is \code{functools.reduce}, which
takes exactly the combiner and null value that exercise asks you to invent.

We should be plain that this is not the same abstraction.  The original's
\code{sum} abstracts over \emph{procedures}; Python's abstracts over
\emph{sequences}.  Both remove the duplication in \code{sum\_integers},
\code{sum\_cubes} and \code{pi\_sum}, but they do it from opposite ends, and
the sequence-shaped answer is the one a Python programmer will reach for.  The
procedure-shaped one is what this chapter is teaching, because it is the one
that generalises to things a sequence cannot express---and because
\code{reduce} and \code{prod}, like \code{pow} in
\link{Section 1.2.6}, are answers with the reasoning removed.  We shall meet
sequences properly in \link{Chapter 2}, and by then we will want them.
\end{revealTheSurface}

\noindent
Once we have \code{sum}, we can use it as a building block in formulating
further concepts.  For instance, the definite integral of a function \( f \)
between the limits \( a \) and \( b \) can be approximated numerically using the
formula

$${\int_a^b \!\!\! f} = {\left[\;f\! \left(a + \frac{dx}{2}\right)
		+ f\! \left(a + dx + \frac{dx}{2}\right)
		+ f\! \left(a + 2dx + \frac{dx}{2}\right) + \,\dots \;\right]\! dx} $$

\noindent
for small values of \( dx \).  We can express this directly as a procedure:

\begin{codeBlock}
>>> def integral(f, a, b, dx):
...     def add_dx(x):
...         return x + dx
...     return sum(f, a + dx / 2.0, add_dx, b) * dx

>>> integral(cube, 0, 1, 0.01)
~\textit{0.24998750000000042}~
\end{codeBlock}

\begin{departure}
The original goes on to halve the step size again, computing
\code{(integral cube 0 1 0.001)} and printing \code{.249999875000001}.  We
cannot: our \code{sum} is a recursive process, one frame deep for each term,
and a thousand terms is more than Python will stack.  The call raises
\code{RecursionError} rather than returning a number.

The remedy is the one the original sets as \link{Exercise 1.30}, which asks
for a \code{sum} that generates an iterative process.  There the exercise is a
matter of good practice; here it is the difference between a procedure that
works at this step size and one that does not.  Written iteratively,
\code{integral(cube, 0, 1, 0.001)} answers
\code{0.24999987500000073}.\footnote{The original's printed value is
\code{.249999875000001}.  Ours differs in the last few digits because the
terms are summed in a different order---one procedure adds them as the
recursion unwinds, from the largest \( x \) down, the other as it goes, from
the smallest up---and floating-point addition is not associative.  Neither
answer is more nearly right than the other; both are approximations to
\( 1/4 \).}
\end{departure}

\noindent
(The exact value of the integral of \code{cube} between 0 and 1 is 1/4.)

\begin{exerciseGroup}
\phantomsection\label{Exercise 1.29}
\exerciseHeading{Exercise 1.29:} Simpson's Rule is a more accurate
method of numerical integration than the method illustrated above.  Using
Simpson's Rule, the integral of a function \( f \) between \( a \) and \( b \) is
approximated as

$$ \frac{h}{3}(y_0 + 4y_1 + 2y_2 + 4y_3 + 2y_4 + \dots + 2y_{n-2} + 4y_{n-1} + y_n), $$

\noindent
where \( h = (b - a) / n \), for some even integer \( n \), and
\( y_k = f(a + kh) \).  (Increasing \( n \) increases the
accuracy of the approximation.)  Define a procedure that takes as arguments
\( f \), \( a \), \( b \), and \( n \) and returns the value of the integral, computed
using Simpson's Rule.  Use your procedure to integrate \code{cube} between 0
and 1 (with \( n = 100 \) and \( n = 1000 \)), and compare the results to those of
the \code{integral} procedure shown above.

\phantomsection\label{Exercise 1.30}
\exerciseHeading{Exercise 1.30:} The \code{sum} procedure above
generates a linear recursion.  The procedure can be rewritten so that the sum
is performed iteratively.  Show how to do this by filling in the missing
expressions in the following definition:

\begin{syntaxForm}
def sum(term, a, next, b):
    def iter(a, result):
        if \meta{??}:
            return \meta{??}
        else:
            return iter(\meta{??}, \meta{??})
    return iter(\meta{??}, \meta{??})
\end{syntaxForm}

\phantomsection\label{Exercise 1.31}
\exerciseHeading{Exercise 1.31:} \begin{enumerate}[a.]
\item
The \code{sum} procedure is only the simplest of a vast number of similar
abstractions that can be captured as higher-order procedures.\footnote{The
intent of \link{Exercise 1.31} through \link{Exercise 1.33} is to demonstrate the
expressive power that is attained by using an appropriate abstraction to
consolidate many seemingly disparate operations.  However, though accumulation
and filtering are elegant ideas, our hands are somewhat tied in using them at
this point since we do not yet have data structures to provide suitable means
of combination for these abstractions.  We will return to these ideas in
\link{Section 2.2.3} when we show how to use \newterm{sequences} as interfaces
for combining filters and accumulators to build even more powerful
abstractions.  We will see there how these methods really come into their own
as a powerful and elegant approach to designing programs.}  Write an analogous
procedure called \code{product} that returns the product of the values of a
function at points over a given range.  Show how to define \code{factorial} in
terms of \code{product}.  Also use \code{product} to compute approximations to
\( \pi \) using the formula\footnote{This formula was discovered by the
seventeenth-century English mathematician John Wallis.}

$$ \frac{\pi}{4}
=
\frac{2\cdot 4\cdot 4\cdot 6\cdot 6\cdot 8\cdots}%
{3\cdot 3\cdot 5\cdot 5\cdot 7\cdot 7\cdots}\,. $$

\item
If your \code{product} procedure generates a recursive process, write one that
generates an iterative process.  If it generates an iterative process, write
one that generates a recursive process.

\end{enumerate}

\phantomsection\label{Exercise 1.32}
\exerciseHeading{Exercise 1.32:} \begin{enumerate}[a.]

\item
Show that \code{sum} and \code{product} (\link{Exercise 1.31}) are both special
cases of a still more general notion called \code{accumulate} that combines a
collection of terms, using some general accumulation function:

\begin{syntaxForm}
accumulate(combiner, null_value, term, a, next, b)
\end{syntaxForm}

\code{accumulate} takes as arguments the same term and range specifications as
\code{sum} and \code{product}, together with a \code{combiner} procedure (of
two arguments) that specifies how the current term is to be combined with the
accumulation of the preceding terms and a \code{null\_value} that specifies what
base value to use when the terms run out.  Write \code{accumulate} and show how
\code{sum} and \code{product} can both be defined as simple calls to
\code{accumulate}.

\item
If your \code{accumulate} procedure generates a recursive process, write one
that generates an iterative process.  If it generates an iterative process,
write one that generates a recursive process.

\end{enumerate}

\phantomsection\label{Exercise 1.33}
\exerciseHeading{Exercise 1.33:} You can obtain an even more
general version of \code{accumulate} (\link{Exercise 1.32}) by introducing the
notion of a \newterm{filter} on the terms to be combined.  That is, combine
only those terms derived from values in the range that satisfy a specified
condition.  The resulting \code{filtered\_accumulate} abstraction takes the same
arguments as accumulate, together with an additional predicate of one argument
that specifies the filter.  Write \code{filtered\_accumulate} as a procedure.
Show how to express the following using \code{filtered\_accumulate}:

\begin{enumerate}[a.]

\item
the sum of the squares of the prime numbers in the interval \( a \) to \( b \)
(assuming that you have a \code{is\_prime} predicate already written)

\item
the product of all the positive integers less than \( n \) that are relatively
prime to \( n \) (i.e., all positive integers \( i < n \) such that
\( \textsc{gcd}(i, n) = 1 \)).

\end{enumerate}
\end{exerciseGroup}

